\mylettrine{E}{s} ist der erste Monat des Jahres des Herrn eintausend dreihundert und dreizehn, da ich, ein Mönch ohne Namen und ohne Ruhm, meine Feder in die Tinte tauche und beginne, diese Blätter mit einer Geschichte zu füllen, die mir seit nun fünf Jahren keine Ruhe lässt. Ich schreibe in einer Klosterzelle, deren Wände von einem langen Winter feucht geworden sind, und der Atem steigt mir sichtbar in die Kälte — ein beständiges Zeichen dafür, dass das Leben vergeht und die Zeit drängt.

Ich bin kein Heiliger und werde es nicht vorgeben. Wer ein Buch mit überschwenglichem Gotteslob eröffnen möchte, dem sei ein anderer Verfasser empfohlen. Ich bin ein lesender, fragender Mensch, der in stillen Stunden lieber in alten Folianten blättert als andächtig die Hände faltet — und der gleichwohl glaubt, dass Gott dem Menschen den Verstand gab, damit er ihn gebrauche. Mein Name tut zur Sache nichts; meine Arbeit möge für sich stehen.

Fünf Jahre lang — fünf lange, oft beschwerliche Jahre — habe ich Zeugen befragt, Urkunden gelesen, Gerüchte abgewogen und Wegstrecken nachgezeichnet, die ein Mann namens Kuonrât von Steinære in seinem unsteten Leben zurücklegte. Was ich dabei fand, hat mich mehr als einmal zum Schweigen gebracht, mehr als einmal zum Kopfschütteln und, ich gestehe es offen, auch mehr als einmal zu einem unfrommen Lachen verleitet. Kuonrâts Geschichte ist wunderlich, zuweilen grotesk, und sie endet, wie die Geschichten solcher Männer zu enden pflegen: mit dem Nachhall einer Lehre, die der Betroffene selbst niemals annehmen wollte.

Doch ehe ich von ihm berichte, muss ich von der Zeit sprechen, aus der er erwuchs. Denn kein Mensch ist vom Himmel gefallen, und kein Charakter entsteht ohne den Boden, der ihn nährt.

Als Kuonrât von Steinære jung war und die Welt noch vor ihm lag wie ein unbeschriebenes Pergament, litt das Reich an einer Wunde, die Jahrzehnte blutete. Es gab keinen rechten Kaiser, keinen rechten König, keinen Hüter der Ordnung — nur Fürsten, die einander befehdeten, Ritter, die raubten, und ein Volk, das litt und schwieg. Jener Zustand, den die Gelehrten das \emph{Interregnum} nennen, war kein bloßer Mangel an einer Krone; er war der Zustand der entfesselten menschlichen Natur ohne das Band der Pflicht, ohne das Gerüst der Ordnung. Die Heilige Schrift spricht davon, wenn es im Buche der Richter heißt, dass in jenen alten Tagen kein König in Israel war, und ein jeder tat, was in seinen eigenen Augen recht schien. Man nehme diese Worte nicht nur für Israel; man nehme sie für jede Zeit, jedes Land, jeden Menschen, der die Pflicht vergisst.

Gott sei gepriesen — und ich tue es, obschon nüchtern —, dass nun, da ich diese Zeilen schreibe, ein Kaiser die Krone trägt. Heinrich, von Gottes Gnaden Kaiser der Römer, wurde im vergangenen Sommer zu Rom gekrönt, und das Reich hat wieder ein Haupt. Die Erleichterung, die mich bei dieser Nachricht überkam, war nicht die fromme Freude eines entrückten Gottesknechtes; sie war die schlichte Erleichterung eines Mannes, der weiß, was es bedeutet, wenn das Gerüst fehlt, das die Ordnung trägt.

Denn die \emph{ordo mundi} — die Ordnung der Welt — ist kein willkürliches Menschenwerk. Die Jahreszeiten folgen festen Gesetzen: der Frühling löst den Winter ab, der Sommer reift das Korn, der Herbst erntet es, und der Winter birgt es in Kälte. Weicht eine Jahreszeit von ihrem Maß ab — kommt der Frost zu früh, bleibt der Regen zu lange aus — so folgen Missernten, Hunger und Tod. Nicht anders ist es mit den Ordnungen des menschlichen Zusammenlebens. Wenn der Bauer nicht sät, hat der Müllergeselle kein Mehl; wenn der Müller nicht mahlt, hat der Bäcker kein Brot; wenn der Ritter nicht schützt, kann der Bauer nicht säen. Der alte Römer Menenius Agrippa erkannte dies, als er dem aufständischen Volk Roms das Bild des Leibes vorhielt: der Magen, so sprach er, nähre alle Glieder, und die Glieder wiederum versorgen den Magen — nehme eines davon seine Pflicht nicht wahr, so sterbe der ganze Leib. Was für den Leib gilt, gilt für das Reich, für das Dorf, für jede Gemeinschaft der Menschen.

Gott gab dem Menschen das \emph{liberum arbitrium} — den freien Willen — und dies ist seine edelste Gabe, denn sie macht den Menschen zum Herrn seiner Entscheidungen. Doch Freiheit ist kein Freibrief. Sie entbindet den Menschen nicht von seiner \emph{triuwe} gegenüber dem Nächsten, von den Pflichten, die ihm sein Stand auferlegt. Der Bauer schuldet dem Herrn seinen Dienst, der Ritter schuldet dem Land seinen Schutz, der Geistliche schuldet der Gemeinde seine Fürsorge. Das höchste Gut ist nicht die Freiheit schlechthin, sondern jene tiefere Freiheit: aus eigenem Erkennen und eigenem Willen die Pflicht zu wählen. So wie einst der Römer Lucius Quinctius Cincinnatus vom Pfluge fortgerufen ward, die Diktatur übernahm, Rom rettete und — das ist das Entscheidende — danach freiwillig zum Pfluge zurückkehrte. Nicht weil man ihn zwang. Weil er es so verstand. Weil in ihm Freiheit und Pflicht keinen Widerspruch bildeten, sondern eine Tugend.

Kuonrât von Steinære verstand es anders. Er nahm die Freiheit und ließ die Pflicht. Er zog das Gold der Treue, die Lust der Beständigkeit, das Abenteuer der Verantwortung vor. Die Zeit des \emph{Interregnum}, in der er lebte und handelte, bot ihm dafür reichlich Raum — mehr Raum, als gut war für ihn und die, die ihm begegneten. Wie ein Fluss, dessen Ufer brechen, schwoll sein Eigenwille ins Maßlose; und wie ein solcher Fluss verwüstete er, was er hätte nähren sollen.

Diese Geschichte ist also eine Mahnung. Keine Heiligenlegende, kein Heldenlied. Aber vielleicht, wenn der Leser sie geduldig zu Ende trägt, weiß er hernach etwas, das er zuvor nur ahnte: dass die Freiheit, die keiner Pflicht gehorcht, am Ende keine Freiheit ist, sondern ein anderer Name für das Verderben.

Möge Gott meiner Feder Geradheit verleihen — und dem Leser einen wachen Sinn.
